\documentclass[../数学分析培优讲义.tex]{subfiles}

\begin{document}

\setcounter{chapter}{7}
\pagestyle{fancy}

\chapter{数项级数}


\section{级数收敛的定义及求和问题}

\begin{theorem}[\markstar\ 加括号]
设 $\displaystyle\sum\limits_{n=0}^{\infty}a_n$ 是一收敛级数,如果把级数的项任意结合而不改变其先后次序,得新级数
\[
(a_1+\cdots+a_{k_1})+(a_{k_1+1}+\cdots+a_{k_2})+\cdots
+(a_{k_{n-1}+1}+\cdots+a_{k_n})+\cdots,
\]
这里的正整数 $k_j$($j=1,2,\ldots$)满足 $k_1<k_2<\cdots$,那么新级数也收敛,且与原级数有相同的和.
\end{theorem}

\begin{theorem}[\markstar]
若上式中同一括号中的项都有相同的符号,那么从加括号后的级数收敛便能推出原级数 $\displaystyle\sum\limits_{n=1}^{\infty}a_n$ 收敛,而且两者有相同的和.
\end{theorem}


\newpage

\subsection{裂项相消法}

\begin{example}
设 $0<x<1$,求如下级数之和:
\[
    \sum\limits_{n=0}^{\infty}\dfrac{x^{2n}}{1-x^{2n+1}}.
\]
\end{example}
\exampleblank

\begin{example}
求下列级数之和:
\begin{enumerate}[label=(\arabic*)]
    \item $\displaystyle\sum\limits_{n=1}^{\infty}\dfrac{1}{(kn-i)(kn+j)}$,其中 $k=i+j$;
    \item $\displaystyle\sum\limits_{n=1}^{\infty}\dfrac{1}{n(n+1)(n+2)}$;
    \item $\displaystyle\sum\limits_{n=1}^{\infty}\dfrac{n-\sqrt{n^2-1}}{\sqrt{n(n+1)}}$;
    \item $\displaystyle\sum\limits_{n=1}^{\infty}(\sqrt{n}-2\sqrt{n+1}+\sqrt{n+2})$.
\end{enumerate}
\end{example}
\exampleblank

\begin{example}
求下列级数 $\displaystyle I=\sum\limits_{n=1}^{\infty}a_n$ 之和:
\begin{enumerate}[label=(\arabic*)]
    \item $\displaystyle a_n=\sin\dfrac{1}{2n}\cos\dfrac{3}{2n}$;
    \item $\displaystyle a_n=\arctan\dfrac{1}{2n^2}$;
    \item $\displaystyle a_n=\arctan\dfrac{2}{4n^2-4n+1}$;
    \item $\displaystyle a_n=3^{n-1}\sin^3\left(\dfrac{\theta}{3^n}\right)$.
\end{enumerate}
\end{example}
\exampleblank

\begin{example}
已知级数 $\displaystyle\sum\limits_{n=1}^{\infty}a_n=A$,$\displaystyle\lim\limits_{n\to\infty}na_n=0$,求级数
\[
    \sum\limits_{n=1}^{\infty}n(a_n-a_{n+1})
\]
之和.
\end{example}
\exampleblank

\begin{example}
设 $a_n>0$ 且 $\displaystyle\sum\limits_{n=1}^{\infty}a_n=S$,若记
\[
    R_n=\sum\limits_{k=n+1}^{\infty}a_k,
\]
则
\[
    I=\sum\limits_{n=1}^{\infty}\dfrac{a_n}{\sqrt{R_{n-1}}+\sqrt{R_n}}=\sqrt{S}.
\]
\end{example}
\exampleblank

\begin{example}
设数列 $\{a_n\}$ 满足
\[
    \sum\limits_{n=1}^{\infty}(a_1+1)(a_2+1)\cdots(a_n+1)=l,
\]
其中 $l>0$ 或 $l=+\infty$.试证明:
\[
    \sum\limits_{n=1}^{\infty}
    \dfrac{a_n}{(a_1+1)(a_2+1)\cdots(a_n+1)}
    =1-\dfrac{1}{l}.
\]
\end{example}
\exampleblank

\begin{example}
设 $a_1>1$,$a_{n+1}=a_n^2-a_n+1$,则
\[
    I=\sum\limits_{n=1}^{\infty}\dfrac{1}{a_n}=\dfrac{1}{a_1-1}.
\]
\end{example}
\exampleblank


\newpage

\subsection{子列方法}

易知,若 $\{S_{2n}\}$ 与 $\{S_{2n+1}\}$ 有相同极限 $S$,则 $\displaystyle\lim\limits_{n\to\infty}S_n=S$.因此若通项 $a_n\to0$,则可由 $S_{2n}\to S$ 得到 $S_{2n+1}\to S$,进而 $\displaystyle\sum\limits_{n=1}^{\infty}a_n=S$.更一般地,若 $a_n\to0$ 且 $\{S_{pn}\}\to S$($p$ 是某个正整数),则 $\displaystyle\sum\limits_{n=1}^{\infty}a_n=S$.我们称这种方法为子列方法.

\begin{example}
计算
\[
    1+\dfrac{1}{2}+\left(\dfrac{1}{3}-1\right)+\dfrac{1}{4}+\dfrac{1}{5}
    +\left(\dfrac{1}{6}-\dfrac{1}{2}\right)+\dfrac{1}{7}+\dfrac{1}{8}
    +\left(\dfrac{1}{9}-\dfrac{1}{3}\right)+\cdots.
\]
\end{example}
\exampleblank

\begin{example}
计算
\[
    1-\dfrac{1}{2}-\dfrac{1}{4}+\dfrac{1}{3}-\dfrac{1}{6}-\dfrac{1}{8}
    +\cdots+\dfrac{1}{2n-1}-\dfrac{1}{4n-2}-\dfrac{1}{4n}+\cdots.
\]
\end{example}
\exampleblank

\begin{example}
将级数 $\displaystyle1-\dfrac{1}{2}+\dfrac{1}{3}-\dfrac{1}{4}+\dfrac{1}{5}-\cdots$ 的各项重新安排,使先依次出现 $p$ 个正项,再出现 $q$ 个负项,然后如此交替.试证新级数的和为
\[
    \ln2+\dfrac{1}{2}\ln\dfrac{p}{q}.
\]
\end{example}
\exampleblank

\begin{example}
证明:若删去调和级数中所有分母含有数字 $9$ 的项,则新级数收敛,且和小于 $80$.
\end{example}
\exampleblank


\newpage

\subsection{发散级数}

\begin{example}
试证明下列级数的发散性:
\begin{enumerate}[label=(\arabic*)]
    \item $\displaystyle\sum\limits_{n=1}^{\infty}(n^2+2)\ln\left(\dfrac{n^2+1}{n^2}\right)$;
    \item $\displaystyle\sum\limits_{n=1}^{\infty}\left(\dfrac{2n^2-3}{2n^2+1}\right)^{n^2}$;
    \item $\displaystyle\sum\limits_{n=1}^{\infty}\cos(n\alpha)$;
    \item $\displaystyle\sum\limits_{n=1}^{\infty}\dfrac{1}{n}$.
\end{enumerate}
\end{example}
\exampleblank

\begin{example}
设 $a_n>0$,$\displaystyle\sum\limits_{n=1}^{\infty}a_n=+\infty$,试证:
\[
    \sum\limits_{n=1}^{\infty}\dfrac{a_n}{S_n}
\]
发散.
\end{example}
\exampleblank

\begin{example}
如果
\[
    \lim\limits_{n\to\infty}a_{n+1}=0,
    \quad
    \lim\limits_{n\to\infty}(a_{n+1}+a_{n+2})=0,
    \quad\ldots,\quad
    \lim\limits_{n\to\infty}(a_{n+1}+\cdots+a_{n+p})=0,
\]
试问 $\displaystyle\sum\limits_{n=1}^{\infty}a_n$ 是否一定收敛.
\end{example}
\exampleblank

\begin{example}
设 $\{a_n\}$ 是递减正数列,且 $\displaystyle\sum\limits_{n=1}^{\infty}a_n=+\infty$,证明:
\[
    \lim\limits_{n\to\infty}
    \dfrac{a_1+a_3+\cdots+a_{2n-1}}{a_2+a_4+\cdots+a_{2n}}=1.
\]
\end{example}
\exampleblank

\begin{example}
证明下列级数发散:
\begin{enumerate}[label=(\arabic*)]
    \item $\displaystyle1-\dfrac{1}{\sqrt2}+\dfrac{1}{3}-\dfrac{1}{\sqrt4}+\dfrac{1}{5}-\dfrac{1}{\sqrt6}+\cdots$;
    \item $\displaystyle1+\dfrac{1}{\sqrt3}-\dfrac{1}{\sqrt2}+\dfrac{1}{\sqrt5}+\dfrac{1}{\sqrt7}-\dfrac{1}{\sqrt4}+\cdots$.
\end{enumerate}
\end{example}
\exampleblank


\newpage

\section{正项级数}

判断级数 $\displaystyle\sum a_n$ 的敛散性,通常有如下方法:
\begin{enumerate}[label=\textcircled{\arabic*}]
    \item 若通项 $a_n\not\to0$,则级数发散;
    \item 考虑部分和 $S_n$ 是否有界,有界则收敛,无界则发散;
    \item Cauchy 收敛准则;
    \item 比较判别法.根据被比较的级数不同,又可细分为 Cauchy 凝聚判别法、Cauchy 积分判别法、Cauchy 根式判别法、d'Alembert 比值判别法、Rabbe 判别法、Gauss 判别法等.
\end{enumerate}


\newpage

\subsection{一般比较判别}


\methodsection{(A) 通项大小比较法}

\begin{theorem}[\marktwostar]
设有两个正项级数
\[
    (A)\sim\sum\limits_{n=1}^{\infty}a_n,
    \qquad
    (B)\sim\sum\limits_{n=1}^{\infty}b_n,
\]
满足 $0<a_n\leq b_n$($n\geq n_0$).则:
\begin{enumerate}[label=(\arabic*)]
    \item 若 $(B)$ 收敛,则 $(A)$ 收敛;
    \item 若 $(A)$ 发散,则 $(B)$ 发散.
\end{enumerate}
\end{theorem}

\begin{remark}
若有 $0<d_n<M$,使得 $0<a_n\leq d_nb_n$,则上述第一条仍成立.
\end{remark}

\begin{corollary}[\marktwostar]
设 $a_n,b_n>0$,且
\[
    \dfrac{a_{n+1}}{a_n}\leq\dfrac{b_{n+1}}{b_n},
\]
则由 $\displaystyle\sum\limits_{n=1}^{\infty}b_n$ 收敛可导出 $\displaystyle\sum\limits_{n=1}^{\infty}a_n$ 收敛.
\end{corollary}

\begin{remark}
因此定理 8.2.1及推论 8.2.1的核心在于利用不等式估计与已知敛散性的级数进而比较,这就要求不等式基础了.
\end{remark}

\begin{example}
判断级数 $\displaystyle I=\sum\limits_{n=3}^{\infty}a_n$ 的敛散性:
\begin{enumerate}[label=(\arabic*)]
    \item $\displaystyle a_n=\dfrac{1}{\sqrt[r]{\ln n}}\quad(r>0)$;
    \item $\displaystyle a_n=\dfrac{1}{(\ln n)^{\ln n}}$;
    \item $\displaystyle a_n=\dfrac{1}{(\ln\ln n)^{\ln n}}$;
    \item $\displaystyle a_n=\dfrac{1}{n^{1+\frac{1}{\ln\ln n}}}$.
\end{enumerate}
\end{example}
\exampleblank

\begin{example}
判断级数 $\displaystyle I=\sum\limits_{n=1}^{\infty}a_n$ 的敛散性:
\begin{enumerate}[label=(\arabic*)]
    \item $\displaystyle a_n=\dfrac{1}{\ln(n!)}\quad(n\geq2)$;
    \item $\displaystyle a_n=\dfrac{\ln(n!)}{n^p}$;
    \item $\displaystyle a_n=\dfrac{\sum\limits_{k=1}^n\ln^2k}{n^p}$;
    \item $\displaystyle a_n=\left[\dfrac{(2n-1)!!}{(2n)!!}\right]^2$.
\end{enumerate}
\end{example}
\exampleblank

\begin{example}
试证明下列命题:
\begin{enumerate}[label=(\arabic*)]
    \item 设 $a_n>0$,若 $\displaystyle\sum\limits_{n=1}^{\infty}a_n$ 收敛,则
    \[
        \sum\limits_{n=1}^{\infty}\dfrac{\sqrt{a_n}}{n^p}
    \]
    在 $p>\dfrac{1}{2}$ 时收敛;
    \item 若 $\displaystyle I=\sum\limits_{n=1}^{\infty}a_n$ 是正项收敛级数,则 $\displaystyle J=\sum\limits_{n=1}^{\infty}\sqrt{a_na_{n+1}}$ 收敛,反之不然.
\end{enumerate}
\end{example}
\exampleblank

\begin{example}
设 $\displaystyle\sum\limits_{n=1}^{\infty}a_n$ 是正项发散级数,试判断下列级数的敛散性:
\begin{enumerate}[label=(\arabic*)]
    \item $\displaystyle\sum\limits_{n=1}^{\infty}\dfrac{a_n}{1+a_n}$;
    \item $\displaystyle\sum\limits_{n=1}^{\infty}\dfrac{a_n}{1+na_n}$;
    \item $\displaystyle\sum\limits_{n=1}^{\infty}\dfrac{a_n}{1+n^2a_n}$;
    \item $\displaystyle\sum\limits_{n=1}^{\infty}\dfrac{a_n}{1+a_n^2}$.
\end{enumerate}
\end{example}
\exampleblank

\begin{example}
设 $a_n\neq0$,且 $a_n\to a\neq0$,则
\[
    \sum\limits_{n=1}^{\infty}|a_{n+1}-a_n|
\]
与
\[
    \sum\limits_{n=1}^{\infty}\left|\dfrac{1}{a_{n+1}}-\dfrac{1}{a_m}\right|
\]
同敛散.
\end{example}
\exampleblank

\begin{example}[\markstar\ Sapagof 判别法]
设正数数列 $\{a_n\}$ 单调递减,则 $\displaystyle\lim\limits_{n\to\infty}a_n=0$ 的充分必要条件是正项级数
\[
    \sum\limits_{n=1}^{\infty}\left(1-\dfrac{a_{n+1}}{a_n}\right)
\]
发散.
\end{example}
\exampleblank

\begin{remark}
由此可得:
\begin{enumerate}[label=\textcircled{\arabic*}]
    \item 设 $\{a_n\}$ 为单调递增的正数数列,则该数列与级数 $\displaystyle\sum\limits_{n=1}^{\infty}\left(1-\dfrac{a_n}{a_{n+1}}\right)$ 同敛散;
    \item 若正项级数 $\displaystyle\sum\limits_{n=1}^{\infty}a_n$ 的部分和数列为 $\{S_n\}$,则 $\displaystyle\sum\limits_{n=1}^{\infty}a_n$ 与 $\displaystyle\sum\limits_{n=1}^{\infty}\dfrac{a_n}{S_n}$ 同敛散.
\end{enumerate}
\end{remark}

\begin{proposition}[\markstar]
\begin{enumerate}[label=(\arabic*)]
    \item 设正项级数 $\displaystyle\sum\limits_{n=1}^{\infty}a_n$ 收敛,记 $\displaystyle R_n=\sum\limits_{k=n}^{\infty}a_k$,则级数 $\displaystyle\sum\limits_{n=1}^{\infty}\dfrac{a_n}{R_n^p}$ 在 $p<1$ 时收敛,$p\geq1$ 时发散;
    \item 设正项级数 $\displaystyle\sum\limits_{n=1}^{\infty}a_n$ 发散,记 $\displaystyle S_n=\sum\limits_{k=1}^na_k$,则级数 $\displaystyle\sum\limits_{n=1}^{\infty}\dfrac{a_n}{S_n^p}$ 在 $p\leq1$ 时发散,$p>1$ 时收敛.
\end{enumerate}
\end{proposition}

\begin{remark}
由证明过程可知,任意正项级数 $\displaystyle\sum\limits_{n=1}^{\infty}a_n$,当 $p>1$ 时,必有 $\displaystyle\sum\limits_{n=1}^{\infty}\dfrac{a_n}{S_n^p}$ 收敛.并且由上述命题,我们可知没有最小的发散级数,也没有最大的发散级数,因此没有所谓的``最完美''的判别法.
\end{remark}

\begin{example}[\markstar]
证明如下问题:
\begin{enumerate}[label=(\arabic*)]
    \item 设正项级数 $\displaystyle\sum\limits_{n=1}^{\infty}a_n$ 发散,则 $\displaystyle\sum\limits_{n=1}^{\infty}\dfrac{a_n}{S_n}$ 发散,$\displaystyle\sum\limits_{n=1}^{\infty}\dfrac{a_n}{S_n^2}$ 收敛;
    \item 设正项级数 $\displaystyle\sum\limits_{n=1}^{\infty}a_n$ 收敛,则 $\displaystyle\sum\limits_{n=1}^{\infty}\dfrac{a_n}{R_n}$ 发散,$\displaystyle\sum\limits_{n=1}^{\infty}\dfrac{a_n}{\sqrt{R_n}}$ 收敛;
    \item 设正项级数 $\displaystyle\sum\limits_{n=1}^{\infty}a_n$ 发散,则存在发散的正项级数 $\displaystyle\sum\limits_{n=1}^{\infty}b_n$,使得 $\displaystyle\sum\limits_{n=1}^{\infty}\dfrac{a_n}{b_n}=0$;
    \item 设正项级数 $\displaystyle\sum\limits_{n=1}^{\infty}a_n$ 收敛,则存在收敛的正项级数 $\displaystyle\sum\limits_{n=1}^{\infty}b_n$,使得 $\displaystyle\sum\limits_{n=1}^{\infty}\dfrac{a_n}{b_n}=0$.
\end{enumerate}
\end{example}
\exampleblank

\begin{example}[\markstar]
举例说明下列命题中的存在性:
\begin{enumerate}[label=(\arabic*)]
    \item 设 $\{b_n\}$ 是趋于 $+\infty$ 的正数列,则必存在收敛的正项级数 $\displaystyle\sum\limits_{n=1}^{\infty}a_n$,使得 $\displaystyle\sum\limits_{n=1}^{\infty}a_nb_n=+\infty$;
    \item 设 $\displaystyle\sum\limits_{n=1}^{\infty}a_n$ 是收敛的正项级数,则存在正数列 $\{\lambda_n\}$ 满足 $\displaystyle\lim\limits_{n\to\infty}\lambda_n=+\infty$ 且 $\displaystyle\sum\limits_{n=1}^{\infty}\lambda_na_n<+\infty$;
    \item 存在 $\{a_n\},\{b_n\}$ 是递减趋于零的数列,且 $\displaystyle\sum\limits_{n=1}^{\infty}a_n,\sum\limits_{n=1}^{\infty}b_n$ 都发散,但 $\displaystyle\sum\limits_{n=1}^{\infty}\min\{a_n,b_n\}$ 收敛;
    \item 设 $\{a_n\}$ 是趋于零的正数列,且 $\displaystyle\sum\limits_{n=1}^{\infty}a_n=+\infty$,则对 $C>0$,存在子列 $\{n_k\}$,使得 $\displaystyle\sum\limits_{k=1}^{\infty}a_{n_k}=C$.
\end{enumerate}
\end{example}
\exampleblank

\begin{example}
设 $0<p_1<p_2<\cdots<p_n<\cdots$,求证:
\[
    \sum\limits_{n=1}^{\infty}\dfrac{1}{p_n}\text{ 收敛}
    \quad\Longleftrightarrow\quad
    \sum\limits_{n=1}^{\infty}\dfrac{n}{p_1+p_2+\cdots+p_n}\text{ 收敛}.
\]
\end{example}
\exampleblank

\begin{remark}
事实上,必要性证明可舍去单调性条件.
\end{remark}

\begin{example}
证明下列不等式:
\begin{enumerate}[label=(\arabic*)]
    \item
    \[
        \sum\limits_{n=1}^m A_n^p
        <\dfrac{p}{p-1}\sum\limits_{n=1}^mA_n^{p-1}a_n,
        \qquad p>1,
    \]
    其中 $\{a_n\}$ 单调递增,$\displaystyle A_n=\dfrac{a_1+a_2+\cdots+a_n}{n}$;
    \item (Hardy-Landan 不等式)设 $p>1$,$\{a_n\}$ 同 (1),则
    \[
        \sum\limits_{n=1}^m\left(\dfrac{a_1+a_2+\cdots+a_n}{n}\right)^p
        <\left(\dfrac{p}{p-1}\right)^p\sum\limits_{n=1}^ma_n^p;
    \]
    \item (Cauleman 不等式)设 $\displaystyle\sum\limits_{n=1}^{\infty}a_n$ 是正项收敛级数,则
    \[
        \sum\limits_{n=1}^{\infty}\sqrt[n]{a_1a_2\cdots a_n}
        \leq e\sum\limits_{n=1}^{\infty}a_n.
    \]
\end{enumerate}
\end{example}
\exampleblank

\begin{example}
研究
\[
    \sum\limits_{n=1}^{\infty}\dfrac{1}{x_n^2}
\]
的敛散性,这里 $x_n$ 是方程 $x=\tan x$ 的正根,并且按递增的顺序编号.
\end{example}
\exampleblank

\begin{example}
设 $\displaystyle\sum\limits_{n=1}^{\infty}a_n$ 为收敛的正项级数,求证:
\[
    \sum\limits_{n=1}^{\infty}a_n^{1-\frac{1}{n}}
\]
收敛.
\end{example}
\exampleblank

\begin{example}
已知 $\varphi(x)$ 是 $(-\infty,+\infty)$ 上周期为 $1$ 的连续函数,满足 $\displaystyle\int_0^1\varphi(x)\,\mathrm{d}x=0$.设
\[
    a_n=\int_0^1e^x\varphi(nx)\,\mathrm{d}x,
\]
求证:级数 $\displaystyle\sum\limits_{n=1}^{\infty}a_n^2$ 收敛.
\end{example}
\exampleblank

\begin{example}[\markstar]
设
\[
    \xi(s)=\sum\limits_{n=1}^{\infty}\dfrac{1}{n^s},
    \qquad s>1,
\]
则
\[
    \xi(s)=s\int_1^{+\infty}\dfrac{[x]}{x^{s+1}}\,\mathrm{d}x.
\]
\end{example}
\exampleblank

\begin{example}[\markstar]
\begin{enumerate}[label=(\arabic*)]
    \item 求证:当 $s>0$ 时,$\displaystyle\int_1^{+\infty}\dfrac{x-[x]}{x^{s+1}}\,\mathrm{d}x$ 收敛;
    \item 求证:当 $s>1$ 时,
    \[
        \int_1^{+\infty}\dfrac{x-[x]}{x^{s+1}}\,\mathrm{d}x
        =\dfrac{1}{s-1}-\dfrac{1}{s}\sum\limits_{n=1}^{\infty}\dfrac{1}{n^s}.
    \]
\end{enumerate}
\end{example}
\exampleblank

\begin{example}
设 $k>0,a>0$,证明:
\begin{enumerate}[label=(\arabic*)]
    \item $\displaystyle\int_a^{+\infty}\dfrac{\sin2n\pi x}{x^k}\,\mathrm{d}x$ 收敛;
    \item $\displaystyle\sum\limits_{n=1}^{\infty}\dfrac{1}{n}\dfrac{\sin2n\pi x}{x^k}$ 收敛.
\end{enumerate}
\end{example}
\exampleblank

\begin{example}
判别
\[
    \sum\limits_{k=1}^{\infty}\left[e-\left(1+\dfrac{1}{1!}+\dfrac{1}{2!}+\cdots+\dfrac{1}{n!}\right)\right]
\]
的敛散性.
\end{example}
\exampleblank


\newpage

\methodsection{(B) 通项趋势比较法}

\begin{theorem}
设 $\displaystyle(A)\sim\sum\limits_{n=1}^{\infty}a_n$,$\displaystyle(B)\sim\sum\limits_{n=1}^{\infty}b_n$ 是两个正项级数,若存在极限
\[
    \lim\limits_{n\to\infty}\dfrac{a_n}{b_n}=l,
\]
则当 $l>0$ 时,级数 $(A)$ 与 $(B)$ 同敛散.
\end{theorem}

\begin{example}
判别级数 $\displaystyle I=\sum\limits_{n=1}^{\infty}a_n$ 的敛散性:
\begin{enumerate}[label=(\arabic*)]
    \item $\displaystyle a_n=\dfrac{1}{n^{1+\frac{1}{n}}}$;
    \item $\displaystyle a_n=\dfrac{\sqrt[n]{2}-1}{n^p}$;
    \item $\displaystyle a_n=\dfrac{\sqrt{n+1}-\sqrt n}{n^p}$;
    \item $\displaystyle a_n=\left(1-\sqrt[3]{\dfrac{n-1}{n+1}}\right)^p\quad(p>0)$;
    \item $\displaystyle a_n=\left[e-\left(1+\dfrac{1}{n}\right)^n\right]^p$.
\end{enumerate}
\end{example}
\exampleblank

\begin{example}
判断级数 $\displaystyle I=\sum\limits_{n=1}^{\infty}a_n$ 的敛散性:
\begin{enumerate}[label=(\arabic*)]
    \item $\displaystyle a_n=\left(\dfrac{1}{n}-\sin\dfrac{1}{n}\right)^p\quad(p>0)$;
    \item $\displaystyle a_n=\left(e^{\frac{1}{n}}-\sin\dfrac{1}{n}\right)^{np}-1$;
    \item $\displaystyle a_n=\ln\left(\dfrac{1}{\cos\dfrac{2\pi}{n}}\right)\quad(n>4)$.
\end{enumerate}
\end{example}
\exampleblank

\begin{example}
设正数列 $\{a_n\}$ 满足
\[
    \lim\limits_{n\to\infty}a_n\sum\limits_{i=1}^na_i^2=1.
\]
试证明:
\[
    \lim\limits_{n\to\infty}\sqrt[3]{3na_n}=1,
\]
且 $\displaystyle\sum\limits_{n=1}^{\infty}a_n^3$ 发散.
\end{example}
\exampleblank


\newpage

\subsection{特殊比较判别}


\methodsection{(A) Cauchy 凝聚判别法}

\begin{theorem}[\markstar\ Cauchy 凝聚判别法]
设 $\{a_n\}$ 是递减正数列,则
\[
    \sum\limits_{n=1}^{\infty}a_n\text{ 收敛}
    \quad\Longleftrightarrow\quad
    \sum\limits_{n=0}^{\infty}2^na_{2^n}\text{ 收敛}.
\]
\end{theorem}

\begin{proposition}[\markstar]
设 $\{a_n\}$ 是递减正数列,若 $\displaystyle\sum\limits_{n=1}^{\infty}a_n$ 收敛,则
\[
    \lim\limits_{n\to\infty}na_n=0.
\]
\end{proposition}

\begin{remark}
反常积分亦有类似结论,甚至证明方法都完全相似.
\end{remark}

\begin{example}
设 $\displaystyle\sum\limits_{n=1}^{\infty}a_n$ 是正项收敛级数,若 $\{na_n\}$ 为单调数列,则
\[
    \lim\limits_{n\to\infty}na_n\ln n=0.
\]
\end{example}
\exampleblank

\begin{example}
判别下列级数的敛散性:
\begin{enumerate}[label=(\arabic*)]
    \item $\displaystyle\sum\limits_{n=1}^{\infty}\dfrac{1}{n^p}$;
    \item $\displaystyle\sum\limits_{n=2}^{\infty}\dfrac{1}{n(\ln n)^p}$;
    \item $\displaystyle\sum\limits_{n=3}^{\infty}\dfrac{1}{n\ln n\ln\ln n}$.
\end{enumerate}
\end{example}
\exampleblank

\begin{example}
试证明下列命题:
\begin{enumerate}[label=(\arabic*)]
    \item 设 $\{a_n\}$ 是递减正数列,若 $\displaystyle\sum\limits_{n=1}^{\infty}a_{2^n}$ 发散,则 $\displaystyle\sum\limits_{n=1}^{\infty}\dfrac{a_n}{n}$ 发散;
    \item 设 $\{a_n\}$ 是递减于零的数列:
    \begin{enumerate}[label=(\roman*)]
        \item 若 $\displaystyle\sum\limits_{n=1}^{\infty}a_n=+\infty$,则 $\displaystyle\sum\limits_{n=1}^{\infty}\min\left\{a_n,\dfrac{1}{n}\right\}=+\infty$;
        \item 若 $\displaystyle\sum\limits_{n=1}^{\infty}\dfrac{a_n}{n}=+\infty$,则 $\displaystyle\sum\limits_{n=1}^{\infty}\dfrac{1}{n}\min\left\{a_n,\dfrac{1}{\ln n}\right\}=+\infty$.
    \end{enumerate}
\end{enumerate}
\end{example}
\exampleblank

\begin{example}[\markstar\ 推广的 Cauchy 凝聚定理]
设 $\{a_n\}$ 是严格递减的正数列,若存在严格递增的正整数列 $\{m_k\}$ 和 $M>0$,使得
\[
    m_{k+1}-m_k\leq M(m_k-m_{k-1}),
    \qquad k=2,3,\ldots,
\]
则级数
\[
    \sum\limits_{n=1}^{\infty}a_n
    \quad\text{与}\quad
    \sum\limits_{k=1}^{\infty}(m_{k+1}-m_k)a_{m_k}
\]
同敛散.
\end{example}
\exampleblank

\begin{remark}
因此若 $\{a_n\}$ 是递减正数列,则 $\displaystyle\sum\limits_{n=1}^{\infty}a_n$ 与下列之一同敛散:
\[
    \sum\limits_{n=1}^{\infty}3^na_{3^n},
    \qquad
    \sum\limits_{n=1}^{\infty}na_{n^2},
    \qquad
    \sum\limits_{n=1}^{\infty}n^2a_{n^3},
    \qquad
    \sum\limits_{n=1}^{\infty}e^na_{e^n}.
\]
\end{remark}

\begin{example}
试证明下列命题:
\begin{enumerate}[label=(\arabic*)]
    \item 设 $\displaystyle\sum\limits_{n=1}^{\infty}a_n$ 是收敛正项级数,记 $\displaystyle R_n=\sum\limits_{k=n+1}^{\infty}a_k$,若 $\displaystyle\sum\limits_{n=1}^{\infty}R_n$ 收敛,则 $\displaystyle\lim\limits_{n\to\infty}na_n=0$;
    \item 设 $\{a_n\}$ 是递减正数列,若 $\displaystyle\sum\limits_{n=1}^{\infty}\dfrac{a_n}{\sqrt n}$ 收敛,则级数 $\displaystyle\sum\limits_{n=1}^{\infty}a_n^2$ 收敛.
\end{enumerate}
\end{example}
\exampleblank

\begin{example}
设 $\{a_n\}$ 递减于零,若 $b_n=a_n-2a_{n+1}+a_{n+2}\geq0$($n\in\mathbb{N}$),则
\[
    \sum\limits_{n=1}^{\infty}nb_n=a_1.
\]
\end{example}
\exampleblank


\newpage

\methodsection{(B) 比式判别法}

\begin{example}
讨论级数 $\displaystyle I=\sum\limits_{n=1}^{\infty}a_n$ 的敛散性:
\begin{enumerate}[label=(\arabic*)]
    \item $\displaystyle a_n=\dfrac{3^n n!}{n^n}$;
    \item $\displaystyle a_n=\dfrac{(2n-1)!!}{3^n n!}$;
    \item $\displaystyle a_n=\dfrac{(n!)^3}{3^{n^4/3}}$;
    \item $\displaystyle a_n=\dfrac{(\alpha+1)(2\alpha+1)\cdots(n\alpha+1)}{(\beta+1)(2\beta+1)\cdots(n\beta+1)}\quad(\alpha>0,\beta>0)$;
    \item $\displaystyle a_n=\dfrac{n^\alpha}{\beta^n}\quad(\beta>0)$.
\end{enumerate}
\end{example}
\exampleblank

\begin{example}
设 $a_1\geq a_2\geq\cdots\geq a_n\geq\cdots>0$,$\displaystyle\lim\limits_{n\to\infty}\dfrac{a_{2n}}{a_n}=\lambda$,则级数 $\displaystyle I=\sum\limits_{n=1}^{\infty}a_n$ 在 $\lambda<\dfrac{1}{2}$ 时收敛,在 $\lambda>\dfrac{1}{2}$ 时发散.
\end{example}
\exampleblank

\begin{example}
\begin{enumerate}[label=(\arabic*)]
    \item 设 $a_n>0$,$\displaystyle\lim\limits_{n\to\infty}\dfrac{a_{2n}}{a_{2n-1}}=\alpha$,$\displaystyle\lim\limits_{n\to\infty}\dfrac{a_{2n+1}}{a_{2n}}=\beta$,若 $\alpha\beta<1$,则 $\displaystyle I=\sum\limits_{n=1}^{\infty}a_n$ 收敛;
    \item 设 $a_n>0$,$\displaystyle\lim\limits_{n\to\infty}\dfrac{a_{n+1}}{a_n}=\lambda<1$,则 $\displaystyle\lim\limits_{n\to\infty}\dfrac{a_{2n}}{a_n}=0$.
\end{enumerate}
\end{example}
\exampleblank

\begin{example}[\markstar\ 比式判别法推广]
设 $a_n>0$,$k_0$ 是固定正整数,若
\[
    \lim\limits_{n\to\infty}\dfrac{a_{n+k_0}}{a_n}=\lambda,
\]
则 $\displaystyle I=\sum\limits_{n=1}^{\infty}a_n$ 在 $\lambda<1$ 时收敛,在 $\lambda>1$ 时发散.
\end{example}
\exampleblank

\begin{example}
设 $\displaystyle I=\sum\limits_{n=1}^{\infty}a_n$ 是正项级数,若
\[
    \lim\limits_{n\to\infty}\left(\dfrac{a_{n+1}}{a_n}\right)^n=\lambda,
\]
则 $I$ 在 $\lambda<e^{-1}$ 时收敛,$\lambda>e^{-1}$ 时发散.
\end{example}
\exampleblank

\begin{example}
判别下列级数的敛散性:
\begin{enumerate}[label=(\arabic*)]
    \item $\displaystyle\dfrac{1}{1+a}+\dfrac{2}{1+a^2}+\dfrac{4}{1+a^4}+\dfrac{8}{1+a^8}+\cdots$;
    \item $\displaystyle\dfrac{a}{1+a}+\dfrac{2a^2}{1+a^2}+\dfrac{4a^4}{1+a^4}+\cdots$.
\end{enumerate}
\end{example}
\exampleblank


\newpage

\methodsection{(C) 根式判别法}

\begin{example}
判别级数 $\displaystyle I=\sum\limits_{n=1}^{\infty}a_n$ 的敛散性:
\begin{enumerate}[label=(\arabic*)]
    \item $\displaystyle a_n=(\sqrt[n]{n}-1)^n$;
    \item $\displaystyle a_n=\left(\dfrac{n}{n+1}\right)^{n(n+1)}$;
    \item $\displaystyle a_n=\left(\dfrac{1+\cos n}{2+\cos n}\right)^{2n-\ln n}$;
    \item $\displaystyle a_n=\dfrac{n^3}{(\sqrt2+(-1)^n)^n3^n}$;
    \item $\displaystyle a_n=\dfrac{\ln(1+\alpha^n)}{n^\beta}\quad(\alpha\geq0)$.
\end{enumerate}
\end{example}
\exampleblank

\begin{example}
设 $\{a_n\},\{b_n\}$ 为正数列,若有
\[
    \lim\limits_{n\to\infty}\dfrac{a_{n+1}}{a_n}=\alpha,
    \qquad
    \lim\limits_{n\to\infty}\sqrt[n]{b_n}=\beta,
\]
试证明级数 $\displaystyle I=\sum\limits_{n=1}^{\infty}a_nb_n$ 在 $\alpha\beta<1$ 时收敛,$\alpha\beta>1$ 时发散.
\end{example}
\exampleblank


\newpage

\methodsection{(D) Cauchy 积分判别法}

\begin{example}
证明下列命题:
\begin{enumerate}[label=(\arabic*)]
    \item 设 $a_1>1$,$\displaystyle a_{n+1}=\dfrac{a_n}{1+a_n^\alpha}$,$0<\alpha<1$,则 $\displaystyle\sum\limits_{n=1}^{\infty}a_n$ 收敛;
    \item 设 $a_n>0$,$\displaystyle\sum\limits_{n=1}^{\infty}\dfrac{a_n}{n}$ 收敛,则
    \[
        \sum\limits_{n=1}^{\infty}\left(\sum\limits_{k=1}^{\infty}\dfrac{a_n}{n^2+k^2}\right)
    \]
    收敛.
\end{enumerate}
\end{example}
\exampleblank

\begin{example}
设
\[
    a_n=2\sqrt n-\sum\limits_{k=1}^n\dfrac{1}{\sqrt k},
\]
则 $\displaystyle\lim\limits_{n\to\infty}a_n=l$ 且 $1\leq l\leq2$.
\end{example}
\exampleblank

\begin{example}
设 $f(x)$ 于 $\lbrack 1,+\infty)$ 上可导,$f'(x)$ 单调递增,且 $f(x)\to A$($x\to+\infty$),证明:
\[
    \sum\limits_{n=1}^{\infty}f'(n)
\]
收敛.
\end{example}
\exampleblank


\newpage

\subsection{其他比较判别 *}

\begin{example}[比值对数判别法]
设有正项级数 $\displaystyle\sum\limits_{n=1}^{\infty}a_n$,且存在
\[
    \lim\limits_{n\to\infty}\left(n\ln\dfrac{a_n}{a_{n+1}}\right)=l,
    \qquad -\infty\leq l\leq+\infty.
\]
若 $l>1$,则级数收敛;若 $l<1$,则级数发散.
\end{example}
\exampleblank

\begin{example}[Rabbe 判别法]
设有正项级数 $\displaystyle I=\sum\limits_{n=1}^{\infty}a_n$ 满足
\[
    \lim\limits_{n\to\infty}n\left(\dfrac{a_n}{a_{n+1}}-1\right)=l,
    \qquad -\infty\leq l\leq+\infty.
\]
若 $l>1$,则 $I$ 收敛;若 $l<1$,则 $I$ 发散.
\end{example}
\exampleblank

\begin{example}[\markstar\ Kummer 判别法]
设 $\displaystyle\sum\limits_{n=1}^{\infty}a_n,\sum\limits_{n=1}^{\infty}b_n$ 是正项级数,$\displaystyle\sum\limits_{n=1}^{\infty}b_n$ 发散,若有
\[
    \lim\limits_{n\to\infty}\left(\dfrac{1}{b_n}\dfrac{a_n}{a_{n+1}}-\dfrac{1}{b_{n+1}}\right)=\lambda,
\]
则 $\displaystyle\sum\limits_{n=1}^{\infty}a_n$ 在 $\lambda>0$ 时收敛,在 $\lambda<0$ 时发散.
\end{example}
\exampleblank

\begin{example}
设
\[
    a_n=C_{\alpha}^{n}=\dfrac{\alpha(\alpha-1)\cdots(\alpha-n+1)}{n!}.
\]
\begin{enumerate}[label=(\arabic*)]
    \item $\alpha$ 在什么条件下,$\displaystyle\lim\limits_{n\to\infty}a_n=0$;
    \item 讨论 $\displaystyle\sum\limits_{n=1}^{\infty}a_n$ 的敛散性.
\end{enumerate}
\end{example}
\exampleblank

\begin{example}[对数判别法]
设 $\displaystyle I=\sum\limits_{n=1}^{\infty}a_n$ 是正项级数,若存在
\[
    \lim\limits_{n\to\infty}\dfrac{\ln\dfrac{1}{a_n}}{\ln n}=l,
\]
则当 $l>1$ 时,$I$ 收敛;当 $l<1$ 时,$I$ 发散.
\end{example}
\exampleblank


\newpage

\section{一般级数}

\begin{proposition}[\marktwostar\ Abel 变换]
设有 $\{a_n\},\{b_n\}$ 两个数列,$\displaystyle B_n=\sum\limits_{k=1}^nb_k$,则:
\begin{enumerate}[label=(\arabic*)]
    \item
    \[
        \sum\limits_{k=1}^na_kb_k
        =a_nB_n+\sum\limits_{k=1}^{n-1}(a_k-a_{k+1})B_k;
    \]
    \item
    \[
        \sum\limits_{k=n+1}^ma_kb_k
        =-A_nb_{n+1}+\sum\limits_{k=n+1}^{m-1}A_k(b_k-b_{k+1})+A_mb_m.
    \]
\end{enumerate}
\end{proposition}

\begin{proposition}[\markstar]
\begin{enumerate}[label=(\arabic*)]
    \item
    \[
        \left|\sum\limits_{k=1}^n\sin kx\right|
        =\left|\sum\limits_{k=1}^n
        \dfrac{\cos\dfrac{2k-1}{2}x-\cos\dfrac{2k+1}{2}x}{2\sin\dfrac{x}{2}}\right|
        \leq\dfrac{1}{|\sin\dfrac{x}{2}|};
    \]
    \item
    \[
        \left|\sum\limits_{k=1}^n\cos kx\right|
        =\left|\sum\limits_{k=1}^n
        \dfrac{\sin\dfrac{2k+1}{2}x-\sin\dfrac{2k-1}{2}x}{2\sin\dfrac{x}{2}}\right|
        \leq\dfrac{1}{|\sin\dfrac{x}{2}|}.
    \]
\end{enumerate}
\end{proposition}

\begin{example}
证明级数
\[
    1-\dfrac{1}{3}\left(1+\dfrac{1}{2}\right)
    +\dfrac{1}{5}\left(1+\dfrac{1}{2}+\dfrac{1}{3}\right)
    -\dfrac{1}{7}\left(1+\dfrac{1}{2}+\dfrac{1}{3}+\dfrac{1}{4}\right)+\cdots
\]
是收敛的.
\end{example}
\exampleblank

\begin{example}[\markstar]
\begin{enumerate}[label=(\arabic*)]
    \item 证明 $\displaystyle\sum\limits_{n=1}^{\infty}\dfrac{(-1)^{[\sqrt n]}}{n}$ 收敛;
    \item 证明 $\displaystyle\sum\limits_{n=1}^{\infty}\dfrac{(-1)^{\sqrt[3]{n}}}{n}$ 收敛.
\end{enumerate}
\end{example}
\exampleblank

\begin{example}
讨论级数
\[
    \sum\limits_{n=1}^{\infty}a_n
    =\dfrac{1}{1^p}-\dfrac{1}{2^q}+\dfrac{1}{3^p}-\dfrac{1}{4^q}+
    \cdots+\dfrac{1}{(2n-1)^p}-\dfrac{1}{(2n)^q}+\cdots,
\]
其中 $p>0,q>0$,的绝对收敛与条件收敛性.
\end{example}
\exampleblank

\begin{example}
研究
\[
    \sum\limits_{n=1}^{\infty}\dfrac{(-1)^{n-1}}{n^{p+\frac{1}{n}}}
\]
的绝对收敛与条件收敛性.
\end{example}
\exampleblank

\begin{example}
设 $f(x)$ 在区间 $\lbrack -1,1\rbrack$ 上有二阶连续导数,且 $\displaystyle\lim\limits_{x\to0}\dfrac{f(x)}{x}=0$,证明:
\[
    \sum\limits_{n=1}^{\infty}f\left(\dfrac{1}{n}\right)
\]
绝对收敛.
\end{example}
\exampleblank

\begin{example}
讨论级数
\[
    \sum\limits_{n=1}^{\infty}\left(1+\dfrac{1}{2}+\dfrac{1}{3}+\cdots+\dfrac{1}{n}\right)\dfrac{\sin nx}{n}
\]
的敛散性.
\end{example}
\exampleblank

\begin{example}
设 $\displaystyle\sum\limits_{n=1}^{\infty}a_n$ 收敛,$\displaystyle\sum\limits_{n=1}^{\infty}(a_n+b_n)$ 绝对收敛,试证:$\displaystyle\sum\limits_{n=1}^{\infty}a_nb_n$ 也收敛.
\end{example}
\exampleblank

\begin{example}
判别下列级数的敛散性:
\begin{enumerate}[label=(\arabic*)]
    \item $\displaystyle\sum\limits_{n=1}^{\infty}\dfrac{(-1)^n}{\sqrt n+(-1)^{n-1}}$;
    \item $\displaystyle\sum\limits_{n=1}^{\infty}\ln\left(1+\dfrac{(-1)^n}{\sqrt[3]{n^2}}\right)$.
\end{enumerate}
\end{example}
\exampleblank

\begin{example}
判断下列级数的敛散性:
\begin{enumerate}[label=(\arabic*)]
    \item $\displaystyle\sum\limits_{n=1}^{\infty}(-1)^n\left(1-\cos\dfrac{\pi}{\sqrt n}\right)$;
    \item $\displaystyle\sum\limits_{n=1}^{\infty}\sin\left(\pi\sqrt{n^2+1}\right)$.
\end{enumerate}
\end{example}
\exampleblank

\begin{example}
判断下列级数的敛散性:
\begin{enumerate}[label=(\arabic*)]
    \item $\displaystyle\sum\limits_{n=2}^{\infty}\dfrac{(-1)^n}{\sqrt n}
    \left[(-1)^n+\sqrt n\sin\dfrac{1}{\sqrt n}\right]$;
    \item $\displaystyle\sum\limits_{n=1}^{\infty}(-1)^{\frac{n(n-1)}{2}}\dfrac{1}{\sqrt n}$.
\end{enumerate}
\end{example}
\exampleblank

记
\[
    a_n^+=\begin{cases}a_n,&a_n\geq0,\\0,&a_n<0,
    \end{cases}
    \qquad
    a_n^-=\begin{cases}-a_n,&a_n\leq0,\\0,&a_n>0,
    \end{cases}
\]
则 $a_n^+,a_n^-\geq0$ 且 $a_n=a_n^+-a_n^-$,$|a_n|=a_n^++a_n^-$.

\begin{theorem}[\markstar]
设 $\displaystyle\sum\limits_{n=1}^{\infty}a_n$ 绝对收敛,则:
\begin{enumerate}[label=(\arabic*)]
    \item $\displaystyle\sum\limits_{n=1}^{\infty}a_n^+$,$\displaystyle\sum\limits_{n=1}^{\infty}a_n^-$ 收敛;
    \item $\displaystyle\sum\limits_{n=1}^{\infty}a_n$ 收敛.
\end{enumerate}
\end{theorem}

\begin{corollary}[\markstar]
若 $\displaystyle\sum\limits_{n=1}^{\infty}a_n$ 条件收敛,则
\[
    \sum\limits_{n=1}^{\infty}a_n^+=+\infty,
    \qquad
    \sum\limits_{n=1}^{\infty}a_n^-=+\infty.
\]
\end{corollary}

\begin{remark}
由此可得条件收敛级数 $\displaystyle\sum\limits_{n=1}^{\infty}a_n$ 收敛的原因在于其中正项与负项大部分抵消了.
\end{remark}

\begin{example}
判别级数 $\displaystyle I=\sum\limits_{n=1}^{\infty}a_n$ 的收敛性:
\begin{enumerate}[label=(\arabic*)]
    \item $\displaystyle a_n=\dfrac{(-1)^n}{\sqrt[3]{n}+(-1)^{n-1}}$;
    \item $\displaystyle a_n=\ln\left(1+\dfrac{(-1)^{n-1}}{(n+1)^\alpha}\right)$;
    \item $\displaystyle a_n=\dfrac{(-1)^n}{[2n+(-1)^n]^\alpha}$;
    \item $\displaystyle a_n=\dfrac{(-1)^{n-1}}{(\sqrt{n+1}+(-1)^n)^\alpha}$.
\end{enumerate}
\end{example}
\exampleblank

\begin{example}
\begin{enumerate}[label=(\arabic*)]
    \item 设 $x\neq k\pi$($k\in\mathbb{Z}$),讨论 $\displaystyle\sum\limits_{n=1}^{\infty}\dfrac{\cos nx}{n^p}$ 与 $\displaystyle\sum\limits_{n=1}^{\infty}\dfrac{\sin nx}{n^p}$ 的条件收敛性与绝对收敛性;
    \item 讨论级数 $\displaystyle\sum\limits_{n=1}^{\infty}\dfrac{(-1)^n\sin n}{n}$ 与 $\displaystyle\sum\limits_{n=1}^{\infty}\dfrac{(-1)^n\sin^2n}{n}$ 的条件收敛性与绝对收敛性.
\end{enumerate}
\end{example}
\exampleblank

关于级数重排的问题,有以下两个结论.

\begin{theorem}[\markstar]
交换绝对收敛级数中无穷多项的次序,所得的新级数仍然绝对收敛,其和也不变.
\end{theorem}

\begin{theorem}[\markstar\ Riemann 定理]
若 $\displaystyle\sum\limits_{n=1}^{\infty}a_n$ 条件收敛,则适当交换各项的次序,可使其收敛到任意事先指定的实数 $S$,也可以使其发散到 $+\infty$ 或 $-\infty$.
\end{theorem}


\newpage

\section{级数的应用 *}


\subsection{与极限的联系}

\begin{example}
设
\[
    x_n=\dfrac{x^n}{n!\,3^n},
\]
求 $\displaystyle\lim\limits_{n\to\infty}x_n$.
\end{example}
\exampleblank

\begin{example}
设 $a_n>0$,$\displaystyle\sum\limits_{n=1}^{\infty}a_n=1$,$\displaystyle A_n=\sum\limits_{k=1}^na_k$,求极限
\[
    \lim\limits_{n\to\infty}\dfrac{e^{A_n}-e^{A_{n-1}}}{A_n^e-A_{n-1}^e}.
\]
\end{example}
\exampleblank

\begin{example}
设数列 $\{a_n\}$ 满足条件:
\begin{enumerate}[label=(\arabic*)]
    \item $0<a_k\leq100a_n$($n=k+1,k+2,\ldots$);
    \item $\displaystyle\sum\limits_{n=1}^{\infty}a_n$ 收敛.
\end{enumerate}
求证:$\displaystyle\lim\limits_{n\to\infty}na_n=0$.
\end{example}
\exampleblank

\begin{example}[\markstar]
设正项级数 $\displaystyle\sum\limits_{n=1}^{\infty}a_n$ 收敛,试证:
\[
    \lim\limits_{n\to\infty}\dfrac{\sum\limits_{k=1}^nka_k}{n}=0.
\]
\end{example}
\exampleblank

\begin{example}
设 $\displaystyle S_n=\sum\limits_{k=1}^n\dfrac{\sin k}{k}=S_n^++S_n^-$,其中 $S_n^+$ 和 $S_n^-$ 分别是正项之和和负项之和.证明:$\displaystyle\lim\limits_{n\to\infty}\dfrac{S_n^+}{S_n^-}$ 存在并求其值.
\end{example}
\exampleblank


\newpage

\subsection{Cesaro 和}

\begin{definition}
设 $\displaystyle\sum\limits_{n=1}^{\infty}a_n$ 为一无穷级数,记
\[
    \sigma_n=\dfrac{S_1+S_2+\cdots+S_n}{n},
\]
其中 $S_n$ 为前 $n$ 项和.若 $\{\sigma_n\}$ 收敛,即 $\displaystyle\lim\limits_{n\to\infty}\sigma_n=\sigma$,则称 $\displaystyle\sum\limits_{n=1}^{\infty}a_n$ 在 Cesaro 意义下收敛,$\sigma$ 称为该级数的 Cesaro 和,记为
\[
    \sum\limits_{n=1}^{\infty}a_n=\sigma(C).
\]
此时称级数可以 Cesaro 求和.
\end{definition}

\begin{remark}
若 $\displaystyle\sum\limits_{n=1}^{\infty}a_n$ 收敛,即 $\displaystyle\lim\limits_{n\to\infty}S_n=S$,则 $\displaystyle\lim\limits_{n\to\infty}\sigma_n=S$.由此可以看出 Cesaro 和为更广义的和,能保证原来收敛的级数仍可求 Cesaro 和,部分发散的级数在 Cesaro 意义下亦可求和,例如
\[
    \sum\limits_{n=1}^{\infty}(-1)^{n-1}=\dfrac{1}{2}(C).
\]
\end{remark}

\begin{example}
求下列级数的 Cesaro 和:
\begin{enumerate}[label=(\arabic*)]
    \item $1+0-1+1+0-1+\cdots$;
    \item $\sin x+\sin2x+\sin3x+\cdots$,其中 $0<x<2\pi$.
\end{enumerate}
\end{example}
\exampleblank

\begin{example}[\markstar\ 可以 Cesaro 求和的必要条件]
设 $\displaystyle\sum\limits_{n=1}^{\infty}a_n$ 可以 Cesaro 求和,则
\[
    \lim\limits_{n\to\infty}\dfrac{S_n}{n}=0,
    \qquad
    \lim\limits_{n\to\infty}\dfrac{a_n}{n}=0.
\]
\end{example}
\exampleblank

\begin{example}[\markstar]
设 $\displaystyle\lim\limits_{n\to\infty}na_n=0$,则
\[
    \sum\limits_{n=1}^{\infty}a_n\text{ 在通常意义下收敛}
    \quad\Longleftrightarrow\quad
    \sum\limits_{n=1}^{\infty}a_n\text{ 可以 Cesaro 求和}.
\]
\end{example}
\exampleblank

\begin{example}[\markstar]
设 $\displaystyle\sum\limits_{n=1}^{\infty}a_n$ 可以 Cesaro 求和,证明:
\[
    \sum\limits_{n=1}^{\infty}a_n\text{ 在通常意义下收敛}
    \quad\Longleftrightarrow\quad
    \lim\limits_{n\to\infty}\dfrac{a_1+2a_2+\cdots+na_n}{n}=0.
\]
\end{example}
\exampleblank


\newpage

\subsection{Cauchy 乘积}

\begin{definition}
级数 $\displaystyle\sum\limits_{n=1}^{\infty}a_n$ 与 $\displaystyle\sum\limits_{n=1}^{\infty}b_n$ 的乘积定义为 $\displaystyle\sum\limits_{n=1}^{\infty}c_n$,其中
\[
    c_n=a_1b_n+a_2b_{n-1}+\cdots+a_nb_1.
\]
这种级数的乘积称为 Cauchy 乘积,又称为对角线乘积.
\end{definition}

\begin{theorem}[\markstar\ Cauchy]
设 $\displaystyle\sum\limits_{n=1}^{\infty}a_n$ 与 $\displaystyle\sum\limits_{n=1}^{\infty}b_n$ 都绝对收敛,它们的和分别为 $A,B$,则将所有 $a_ib_j$ 按任何方式求和而成的级数也绝对收敛,且和为 $AB$.
\end{theorem}

\begin{theorem}[\markstar\ Mertens]
设 $\displaystyle\sum\limits_{n=1}^{\infty}a_n$ 与 $\displaystyle\sum\limits_{n=1}^{\infty}b_n$ 收敛且至少有一个绝对收敛,它们的和分别为 $A,B$,则它们的 Cauchy 乘积 $\displaystyle\sum\limits_{n=1}^{\infty}c_n$ 收敛,且和为 $AB$.
\end{theorem}

\begin{theorem}[\markstar\ Abel]
设 $\displaystyle\sum\limits_{n=1}^{\infty}a_n$,$\displaystyle\sum\limits_{n=1}^{\infty}b_n$ 与它们的 Cauchy 乘积 $\displaystyle\sum\limits_{n=1}^{\infty}c_n$ 都收敛,且和分别为 $A,B,C$,则 $C=AB$.
\end{theorem}

\begin{example}[\markstar]
设 $|x|<1$,则
\[
    \dfrac{1}{1-x}\sum\limits_{n=0}^{\infty}a_nx^n
    =\left(\sum\limits_{n=0}^{\infty}x^n\right)
    \left(\sum\limits_{n=0}^{\infty}a_nx^n\right)
    =\sum\limits_{n=0}^{\infty}(a_0+a_1+\cdots+a_n)x^n.
\]
特别地,
\[
    \left(\dfrac{1}{1-x}\right)^2
    =\dfrac{1}{1-x}\sum\limits_{n=0}^{\infty}x^n
    =\sum\limits_{n=0}^{\infty}(n+1)x^n.
\]
\end{example}
\exampleblank

\begin{example}
证明:$\displaystyle\sum\limits_{n=1}^{\infty}\dfrac{a^n}{n!}$ 与 $\displaystyle\sum\limits_{n=1}^{\infty}\dfrac{b^n}{n!}$ 绝对收敛,且它们的 Cauchy 乘积等于
\[
    \sum\limits_{n=0}^{\infty}\dfrac{(a+b)^n}{n!}.
\]
\end{example}
\exampleblank

\begin{example}
求 $\displaystyle\sum\limits_{n=1}^{\infty}nx^{n-1}$ 与 $\displaystyle\sum\limits_{n=1}^{\infty}(-1)^{n-1}nx^{n-1}$ 的 Cauchy 乘积.
\end{example}
\exampleblank

\begin{example}
求下列级数 $S=\sum c_n$ 的和:
\begin{enumerate}[label=(\arabic*)]
    \item $\displaystyle\sum\limits_{n=0}^{\infty}c_n$,其中 $\displaystyle c_n=\sum\limits_{k=0}^nx^ky^{n-k}$,$|x|<1,|y|<1$;
    \item $\displaystyle\sum\limits_{n=1}^{\infty}c_n$,其中 $\displaystyle c_n=\sum\limits_{k=1}^{\infty}\dfrac{1}{k(k+1)(n-k+1)!}$.
\end{enumerate}
\end{example}
\exampleblank


\newpage

\subsection{无穷乘积}

\begin{definition}
设 $\displaystyle\prod\limits_{n=1}^{\infty}p_n$ 是一个无穷乘积,称
\[
    P_n=\prod\limits_{k=1}^np_k=p_1\cdots p_n
\]
为这个无穷乘积的部分乘积.若 $\displaystyle\lim\limits_{n\to\infty}P_n=p$(有限且 $\neq0$),则称这个无穷乘积是收敛的,反之称它是发散的.
\end{definition}

为什么要去掉 $p=0$ 的情况呢?实则是为了得到类似级数收敛的必要条件.

\begin{theorem}[\markstar]
无穷乘积 $\displaystyle\prod\limits_{n=1}^{\infty}p_n$ 收敛的必要条件是
\[
    \lim\limits_{n\to\infty}p_n=1.
\]
\end{theorem}

\begin{theorem}[\markstar]
\[
    \prod\limits_{n=1}^{\infty}(1+a_n)\text{ 收敛}
    \quad\Longleftrightarrow\quad
    \sum\limits_{n=1}^{\infty}\ln(1+a_n)\text{ 收敛}.
\]
\end{theorem}

\begin{remark}
收敛情况下,若后者和为 $S$,则前者积为 $e^S$.
\end{remark}

\begin{theorem}[\markstar]
设 $a_n>0$($n>n_0$),则 $\displaystyle\prod\limits_{n=1}^{\infty}(1+a_n)$ 与 $\displaystyle\sum\limits_{n=1}^{\infty}a_n$ 同敛散.
\end{theorem}

\begin{theorem}[\markstar]
如果 $\displaystyle\sum\limits_{n=1}^{\infty}a_n^2<+\infty$,那么 $\displaystyle\prod\limits_{n=1}^{\infty}(1+a_n)$ 与 $\displaystyle\sum\limits_{n=1}^{\infty}a_n$ 同敛散.
\end{theorem}

\begin{theorem}
设 $-1<a_n<0$,且 $\displaystyle\sum\limits_{n=1}^{\infty}a_n$ 发散,则 $\displaystyle\prod\limits_{n=1}^{\infty}(1+a_n)$ 发散到 $0$.
\end{theorem}

\begin{theorem}
如果 $\displaystyle\sum\limits_{n=1}^{\infty}a_n$ 收敛,设 $\displaystyle\sum\limits_{n=1}^{\infty}a_n^2$ 发散,那么 $\displaystyle\prod\limits_{n=1}^{\infty}(1+a_n)$ 发散到 $0$.
\end{theorem}

\begin{example}
证明下列等式:
\begin{enumerate}[label=(\arabic*)]
    \item $\displaystyle\prod\limits_{n=2}^{\infty}\dfrac{n^3-1}{n^3+1}=\dfrac{2}{3}$;
    \item $\displaystyle\prod\limits_{n=0}^{\infty}\left[1+\left(\dfrac{1}{2}\right)^{2^n}\right]=2$.
\end{enumerate}
\end{example}
\exampleblank

\begin{example}
讨论下列无穷乘积的敛散性:
\begin{enumerate}[label=(\arabic*)]
    \item $\displaystyle\prod\limits_{n=1}^{\infty}\dfrac{1}{n}$;
    \item $\displaystyle\prod\limits_{n=1}^{\infty}\dfrac{(n+1)^2}{n(n+2)}$;
    \item $\displaystyle\prod\limits_{n=1}^{\infty}\sqrt[n]{1+\dfrac{1}{n}}$.
\end{enumerate}
\end{example}
\exampleblank

\begin{example}
讨论下列无穷乘积的敛散性:
\begin{enumerate}[label=(\arabic*)]
    \item $\displaystyle\prod\limits_{n=1}^{\infty}\dfrac{n}{\sqrt{n^2+1}}$;
    \item $\displaystyle\prod\limits_{n=2}^{\infty}\left(\dfrac{n^2-1}{n^2+1}\right)^p\quad(p\in\mathbb{R})$.
\end{enumerate}
\end{example}
\exampleblank

\end{document}
